% =====================================================================
% DISCUSSION --- DRAFT (regenerated from results/findings.md,
% audit/METHOD_VERIFICATION_FULL.md, audit/CLASS_D_EXCLUSIONS.md; 2026-08).
% Combined TSB-StreamingAD-U + TSB-StreamingAD-M. Compile-safe fragment:
% no preamble; \input into main.tex in place of \section{Discussion}.
% Requires siunitx, natbib, cleveref. Every numeral traces to findings.md.
% =====================================================================

\section{Discussion}\label{sec:discussion}

The benchmark yields a constructive result, not merely a cautionary one. It
identifies a single detector family that works, explains why it works, quantifies
how cheaply it works, and turns the fidelity audit that made the ranking
trustworthy into a reusable methodology. We open with the takeaways, then develop
the positive findings before analysing, in depth, why the confidence-based family
fails.

\subsection{Key findings and takeaways}\label{sec:disc-takeaways}
\begin{itemize}
  \item \textbf{There is a clear winner.} On the combined \num{527}-dataset
  benchmark, detectors that model the in-distribution feature manifold occupy the
  entire top tier; Mahalanobis distance (\num{0.826}) and DFM-PCA (\num{0.812})
  lead decisively across both univariate and multivariate families.
  \item \textbf{A concrete recommendation.} For streaming time-series OOD
  detection, \emph{model the geometry and density of the in-distribution feature
  manifold} (Mahalanobis or DFM-PCA) rather than trusting the classifier's
  confidence; pair the detector with global normalisation for amplitude/level
  shifts.
  \item \textbf{The winners are cheap and self-contained.} They attain top accuracy
  at $\sim$\SIrange{1}{4}{\milli\second} per window (Mahalanobis
  \SI{1.9}{\milli\second}, DEEDEE \SI{1.0}{\milli\second}), require no auxiliary
  outlier data and no retraining, and are Pareto-optimal in accuracy versus latency.
  \item \textbf{Dependable across regimes.} The feature-manifold family is strong on
  the STABLE specificity controls (Mahalanobis \num{0.855}) and the subtle DRIFT
  regime (\num{0.858}), so it neither misses drift nor over-flags near-stationary
  streams.
  \item \textbf{The audit is a first-class contribution.} CatSight, InvAD, DEEDEE,
  and even Mahalanobis reach the top tier \emph{only} after principled,
  paper-checked corrections; the verify-before-experiment gate is a transferable
  methodology that most baseline studies omit.
  \item \textbf{A streaming extension of prior work.} These results independently
  confirm and extend the recommendation of TS-OOD \citep{Gungor2025} from the
  static, same-dataset semantic-shift setting to streaming, window-level
  distributional shift, across \num{527} datasets and \num{17} detectors.
  \item \textbf{A disciplined cautionary result.} Every post-hoc softmax, logit,
  activation, and gradient detector falls \emph{below chance} with systematically
  inverted scores; this is a genuine, audited failure of the confidence premise,
  not an implementation artifact.
\end{itemize}

\subsection{What works, and why: the feature-manifold family}\label{sec:disc-why}
The central positive result is that detectors which model the in-distribution
\emph{feature} manifold dominate, and the mechanism is clean and general. These
methods characterise the geometry and density of ID pre-logit features directly ---
Mahalanobis by a tied-covariance class-conditional Gaussian, DFM-PCA by
per-class principal-component reconstruction error, DiMMAD by a continuous-metric
distance ensemble --- and score a test window by how far it lies from that learned
manifold. Because they never consult the overconfident output head, they are
structurally immune to the overconfidence pathology that sinks the post-hoc family
(\cref{sec:disc-fail}): even when the classifier saturates on a far-out window, the
feature vector still records the distributional gap, and distance/density in that
space grows monotonically with the shift. The normalisation dichotomy of
\cref{fig:norm} is the direct fingerprint of this mechanism: under global
normalisation the second source retains its amplitude and level gap, so S2 windows
sit many standard deviations outside the ID feature cloud, and the feature methods
separate them cleanly (Mahalanobis \num{0.867}, DFM-PCA \num{0.863} global). The
practical virtues compound the accuracy advantage. The winners need no auxiliary
outlier corpus and no retraining --- they are fit on ID features alone --- and they
are among the cheapest detectors in the benchmark (\cref{sec:results-efficiency}),
so the accuracy--latency tradeoff that usually forces a compromise simply does not
arise. For a practitioner the recommendation is therefore unambiguous and
low-risk: default to Mahalanobis distance or DFM-PCA.

\subsection{Category-level strengths as deployment guidance}\label{sec:disc-category}
The per-category breakdown (\cref{fig:heatmap}) converts the winning family's
robustness into actionable guidance. The feature-manifold detectors are strongest
precisely where deployment demands reliability: on the STABLE specificity controls,
where a well-calibrated detector must \emph{not} raise alarms, Mahalanobis
(\num{0.855}), DFM-PCA (\num{0.849}), and DiMMAD (\num{0.801}) all maintain strong
false-positive control; and on the subtle near-distribution DRIFT regime, where a
gradual shift is easiest to miss, Mahalanobis reaches \num{0.858} and DFM-PCA
\num{0.825}. The strong-shift OOD category is the family's (relatively) weakest,
around \num{0.76} for the two leaders --- still comfortably above chance, but a
useful reminder that a cleanly different S2 regime can push representations into a
less-calibrated region of the tied-covariance model. The deployment reading is
favourable and specific: use the feature-manifold family as a default; expect its
highest confidence on stable and drifting streams; and, where the deployment is
dominated by abrupt regime changes, consider ensembling the tied-covariance
Mahalanobis with the reconstruction-based DFM-PCA, whose principal-subspace view is
complementary.

\subsection{The fidelity audit and its corrections as a methodology}\label{sec:disc-audit}
A defining methodological contribution of this thesis is that implementation
faithfulness is treated as a \emph{measured precondition} rather than an
assumption. We audited every candidate detector against its source paper and,
where a repository exists, its official code (\cref{app:fidelity}). Of the
\num{24} audited candidates, \num{17} were faithful or admitted a validated,
paper-faithful correction and entered the benchmark; the remaining \num{7} were
protocol-incompatible and excluded (\cref{tab:classd}). This is not bookkeeping:
several of the top-tier results exist \emph{only because} of the audit. CatSight
(\num{0.735}) reaches the top tier after correcting an asserted orientation on its
common-spatial-pattern distance; InvAD (\num{0.740}) is competitive once its
affine-coupling invertibility is recognised to reduce the score to a feature-space
statistic; DEEDEE (\num{0.700}) required its adjacent-dimension variant to be
labelled and validated; and even Mahalanobis reaches the sole top position only
after its within-class-scatter (tied-covariance) estimation was corrected. Most
baseline studies never audit their baselines; by converting faithfulness from an
unstated assumption into a per-method verdict with a synthetic-validation gate, we
obtain a ranking whose top entries are trustworthy and whose failures are
attributable. The audit is thus a reusable protocol --- verify, correct where a
paper-faithful fix exists, exclude where it does not, and report the ablation ---
as much as a set of results.

\subsection{Honest corrections to the prior write-up}\label{sec:disc-corrections}
Reporting how the audit changed the headline is itself a result; four corrections
stand out, each traceable to a specific per-method fix (\cref{app:fidelity}).
\paragraph{(a) Mahalanobis is the sole top method.}
An earlier write-up reported a tie at the top between Mahalanobis and DriftLens.
The audit established that the per-sample DriftLens implementation is a
PCA-whitened Mahalanobis detector correlating at $\rho \approx 0.999$ with
Mahalanobis proper --- the same method under another name --- and that the official
DriftLens score is batch/window-level Fr\'echet drift, which the shuffled
per-window protocol cannot realise. With DriftLens excluded as
protocol-incompatible and the tied-covariance estimation corrected, Mahalanobis
(\num{0.826}) stands alone ahead of DFM-PCA (\num{0.812}); the earlier tie was a
relabelling artifact, not two independent methods converging.
\paragraph{(b) SRS falls below chance once restored to its true form.}
A prior write-up promoted SRS as the faithful, mechanistically distinct
time-series method and reported it near the top. Restored to its true
seasonal-ratio form --- STL decomposition plus conditional-VAE scoring on the raw
series --- SRS falls to \num{0.352}, below the chance line and behind even DiffAD
(\num{0.531}). The seasonal-ratio assumption presumes a stable seasonal template
learned from the ID regime, which the abrupt two-source boundary violates; the one
genuinely distinct time-series-specific method does not generalise to streaming
shift. This is the most consequential reversal, and we report it plainly.
\paragraph{(c) CatSight, InvAD, and DEEDEE rank high only after fixes.}
These three now sit in the top tier (\numlist{0.735;0.740;0.700}) but owe their
standing to orientation, reduction, and variant corrections rather than to the
mechanisms their names advertise; each is labelled an adaptation in
\cref{app:fidelity}. Reporting them under their original names without these caveats
would overstate the fidelity of the result, so we do not.
\paragraph{(d) Faithful DICE (energy) fails with the post-hoc family.}
The corrected, paper-faithful DICE --- signed top-$k$ contribution sum followed by
an energy score, matching \citet{Sun2022} --- does not escape the collapse: it
scores \num{0.313}, squarely inside the inverted cluster alongside MSP
(\num{0.370}), SCALE (\num{0.308}), ODIN (\num{0.303}), energy (\num{0.301}), and
GradNorm (\num{0.294}). This is the crucial control on the audit: the post-hoc
failure is \emph{not} an artifact of buggy code. Even repaired to match its paper,
the family inverts, because its defining assumption breaks under streaming shift.

\subsection{Why post-hoc confidence fails below chance}\label{sec:disc-fail}
The failure of the confidence-based family is deep enough to warrant a dedicated
analysis, because ``below chance'' is a far stronger statement than ``ineffective.''
We develop seven interlocking mechanisms.

\paragraph{1. Below chance means systematically inverted, not random.}
A detector at AUROC $\approx 0.5$ is uninformative; the post-hoc cluster sits at
\numrange{0.294}{0.370}, well \emph{below} chance. By the probabilistic reading of
AUROC, an in-distribution window receives a \emph{higher} OOD score than a genuinely
shifted window the majority of the time. These detectors are therefore reliably
\emph{wrong}, not noisy --- there is real, exploitable signal, but it points the
wrong way. A detector that is consistently inverted is more dangerous in deployment
than one that is merely uninformative, because its confident answers are
anti-correlated with the truth.

\paragraph{2. The overconfidence dichotomy.}
The inversion has a well-documented root: deep networks are frequently \emph{more}
confident on shifted or far-out inputs than on in-distribution ones. A deep ReLU
network is provably free to be arbitrarily confident far from the training manifold
\citep{Hein2019, Nguyen2015}. When an S2 window arrives many standard deviations
outside the ID region, the network can emit a near-maximal softmax probability;
MSP, ODIN, energy, and ReAct then read that saturated confidence as strong evidence
of \emph{in}-distribution, assigning the OOD window a low anomaly score and the
genuinely ID window a comparatively higher one. The confidence signal is not weak;
it is pointed backwards.

\paragraph{3. The backbone is trained on temporal pseudo-classes.}
The pathology is sharpened by how the backbone is trained. There are no semantic
labels in a streaming series, so the backbone is trained with cross-entropy over
$K$ \emph{temporal pseudo-classes} formed by binning the ID training windows into
arbitrary contiguous segments. The resulting confidence measures which temporal
bin a window resembles, a quantity with no principled relationship to true
distributional novelty. An OOD window can be assigned confidently to some
pseudo-class simply because it happens to align with one arbitrary bin; confidence
is decoupled from, and empirically anti-correlated with, genuine novelty.

\paragraph{4. Distribution shift is not semantic novelty.}
The post-hoc methods were designed for image classification, where OOD means a
\emph{new semantic class} and low confidence is a reasonable proxy for ``none of my
classes.'' Streaming distributional shift is a covariate shift within a
label-free stream: there is no semantic class to be uncertain about. The premise
that low confidence signals novelty simply does not transfer to this setting. Our
result confirms and extends \citet{Gungor2025}, who found post-hoc scores weak
under static semantic shift, to the harder streaming regime where they do not merely
weaken but invert.

\paragraph{5. A shared broken signal explains the tight cluster.}
MSP, ODIN, energy, ReAct, SCALE, DICE, and GradNorm all read the same
over-confident logits (or their gradients), differing only in post-processing. Since
the underlying signal is the same inverted confidence, they fail \emph{together} and
cluster within a narrow \num{0.08}-wide band (\numrange{0.294}{0.370}). This
co-location is itself evidence for the mechanism: independent weaknesses would
scatter, whereas a shared broken input produces a tight, correlated failure. By the
same logic the feature-manifold family, which reads a different and intact signal,
separates cleanly and wins together.

\paragraph{6. Global normalisation amplifies the inversion.}
The normalisation dichotomy (\cref{fig:norm}) shows the mechanism operating as a
dial. Under global normalisation the post-hoc cluster is most inverted (energy
\num{0.214}, GradNorm \num{0.213}, ODIN \num{0.218}), because the preserved
out-of-range magnitudes push the logits to their most saturated, most overconfident
extreme; under per-series normalisation, which rescales away the level gap, the same
detectors recover toward --- but never reach --- chance (energy \num{0.452},
GradNorm \num{0.433}). The very preprocessing that most helps the feature methods
most hurts the confidence methods, because a larger level gap is simultaneously
cleaner geometry and stronger overconfidence.

\paragraph{7. This is not an orientation bug, and we do not flip the sign.}
It would be tempting to ``fix'' an inverted detector by negating its score. We do
not, and the reason is a matter of scientific integrity. The Pass-1 audit verified
that MSP and ODIN are faithful implementations with the \emph{correct},
paper-defined orientation (\cref{app:fidelity}); their inversion is the method's own
behaviour, not a coding error. Negating the score to recover AUROC would be
post-hoc orientation-fitting --- choosing the sign after seeing the labels --- which
is precisely the malpractice the audit was built to catch (we document its danger in
the regime-dependent DiverseMix correction, whose principled orientation reverses
between synthetic and real data). The honest conclusion is that the method's own
defined orientation genuinely fails on streaming shift; reporting the true,
below-chance number is the point, not a problem to be papered over.

\subsection{A single winning principle}\label{sec:disc-principle}
A publication-relevant subtlety is that several detectors marketed as distinct ---
DFM-PCA (PCA reconstruction), DiMMAD (distance ensemble), the excluded DriftLens
(PCA-Mahalanobis) and TD-IVDM (KDE density), and Mahalanobis itself --- are, in
their per-sample implementations, all variants of distance or density estimation in
the frozen feature space. The impression that ``a diverse set of methods wins'' is
therefore partly an artifact of relabelling: the robust signal is feature-space
distance/density, surfacing under several names. This sharpens rather than weakens
the headline --- it identifies a single winning principle rather than a coincidence
of unrelated methods, and it is the reason two of the seven exclusions
(\cref{tab:classd}) are protocol-incompatible duplicates rather than genuinely new
detectors. For the practitioner it is reassuring: the recommendation reduces to one
idea, robustly instantiated.

\subsection{Extending TS-OOD to the streaming regime}\label{sec:disc-extension}
The study sits in the empirical-benchmark lineage of \citet{Gungor2025} and extends
it along three axes. Where TS-OOD evaluated eight detectors on thirteen datasets
under a static, same-dataset semantic-shift protocol, we evaluate \num{17} audited
detectors on \num{527} datasets under a leakage-free, window-level streaming-shift
protocol spanning univariate and multivariate families. The core recommendation ---
that deep feature modeling is the most promising direction for time-series OOD
detection --- not only survives this much harder setting but is stated more
strongly: under streaming shift the post-hoc alternatives do not merely lag, they
invert. The breadth of coverage (\num{527} datasets, both modalities, three
difficulty categories, two normalisation schemes) and the strongly significant
Friedman result together make this, to our knowledge, the most comprehensive fair
evaluation of OOD detectors on streaming time series to date, and a positive,
constructive corroboration of the feature-modeling thesis.

\subsection{Threats to validity and limitations}\label{sec:disc-limits}
We state the scope honestly; these are scoping decisions that define the natural
next experiments rather than undermining the present conclusions, which rest on the
very strongly significant Friedman result and the consistent rank ordering across
both families, categories, and normalisation schemes.
\begin{enumerate}
  \item \textbf{CPU-only.} All experiments run on CPU; the inference-time comparisons
  (\cref{fig:efficiency}) are CPU latencies and would change in absolute terms,
  though not in ranking, on accelerated hardware.
  \item \textbf{Friedman complete-case caveat.} The \num{17}-detector complete-case
  Friedman pool is $N = \num{250}$ univariate datasets, because SRS is
  univariate-seasonal and is not run on the multivariate split. A combined-family test
  over the full corpus (\num{16} detectors excluding SRS, $N = \num{515}$ univariate
  and multivariate) is equally highly significant ($\chi^2 = \num{2336.59}$,
  $p < \num{1e-16}$), so the conclusion holds across both splits; only an
  SRS-inclusive combined test is out of scope.
  \item \textbf{Small Source-1 training windows.} The source-boundary split trains
  only on the pre-boundary Source-1 region, typically a small leading fraction of
  each file; the resulting ID pool is modest, which particularly constrains density
  estimators fit on few windows.
  \item \textbf{Class-D methods run under separate protocols.} The seven excluded
  detectors (\cref{tab:classd}) are reported in \cref{tab:class_d_appendix} under
  their own protocol-appropriate arms on both TSB-U and -M; even so, none rises
  clearly above chance (mean AUROC \numrange{0.494}{0.696}). Their behaviour under a
  protocol they \emph{are} designed for (ordered streams, auxiliary outlier corpora,
  batch-level drift) remains untested here. Reporting them separately is the honest
  choice, but it means the comparison to the headline set is indicative, not
  like-for-like.
  \item \textbf{Single seed.} All splits and initialisations use seed \num{42}; a
  multi-seed study would attach confidence intervals to every AUROC.
\end{enumerate}

\subsection{Promise and future directions}\label{sec:disc-future}
The limitations map directly onto concrete next steps. Completing the SRS
multivariate run would close the Friedman complete-case gap and yield a definitive
combined-family significance test. A multi-seed study would replace point estimates
with confidence intervals. Re-training the backbone with a multi-positive
contrastive loss is expected, per \citet{Gungor2025}, to lift the activation-based
detectors specifically, testing whether the post-hoc collapse is intrinsic to the
confidence premise or an artifact of the pseudo-class backbone. An
orientation-stability study would formalise when energy-based scores can be trusted
without post-hoc sign-fitting. Finally, an online deployment of the Pareto-optimal
feature-manifold detectors --- Mahalanobis and DFM-PCA with incremental covariance
updates --- would test their promise under true streaming conditions. For
practitioners deploying OOD detection on streaming time series today, the
recommendation is already concrete and positive: \emph{default to Mahalanobis
distance or DFM-PCA}, prefer global normalisation for amplitude/level shifts, and do
not rely on post-hoc softmax, logit, or gradient scores, which invert under exactly
the far-out inputs they are meant to flag.
